
\section{Introduction}
\label{sec:introduction}

\noindent\textbf{Purpose.}
This document describes the software architecture of the
\emph{Medication Safety Assistant}, including its main components,
responsibilities, interfaces, and relationships. It serves as
a technical reference for the project team and course reviewers.

\medskip
\noindent\textbf{Scope.}
The Medication Safety Assistant is a Web-based research prototype
supporting medication intake, identification, safety assessment,
simulated healthcare professional review, and medication reminders.

This document covers architectural goals and constraints,
the use-case model, and the logical view. The system provides
decision support only and does not replace professional medical
judgment. Deployment and implementation views are deferred
as permitted by the PA3 requirements.

\medskip
\noindent\textbf{Definitions, Acronyms and Abbreviations.}

\noindent
\begin{tabularx}{\textwidth}{|L{2.5cm}|Y|}
    \hline
    \textbf{Term} & \textbf{Definition} \\ \hline
    HCP & Healthcare Professional \\ \hline
    DDI & Drug--Drug Interaction \\ \hline
    OCR & Optical Character Recognition \\ \hline
    NLP & Natural Language Processing \\ \hline
    RBAC & Role-Based Access Control \\ \hline
\end{tabularx}

\medskip
\noindent\textbf{References.}

\noindent
\begin{tabularx}{\textwidth}{|L{1.5cm}|Y|L{2.8cm}|}
    \hline
    \textbf{ID} & \textbf{Document} & \textbf{Version / Date} \\ \hline
    R1 & Vision Document (PA2-Group02-VisionDocument)
       & 1.1 / 01 Nov 2026 \\ \hline
    R2 & Use-Case Specification (PA2-Group02-UseCaseSpecification)
       & 1.0 / 01 Nov 2026 \\ \hline
    R3 & Software Development Plan (PA2-Group02-ProjectPlan)
       & 1.1 / 01 Nov 2026 \\ \hline
\end{tabularx}

\medskip
\noindent\textbf{Overview.}
Section 2 presents architectural goals and constraints.
Section 3 presents the use-case model.
Section 4 describes the logical architecture and its components.
Sections 5 and 6 are reserved for future deployment and
implementation details.
